\documentclass[11pt]{article}

\usepackage[a4paper,margin=1in]{geometry}
\usepackage{amsmath,amssymb,mathtools}
\usepackage{booktabs}
\usepackage{enumitem}
\usepackage[hidelinks]{hyperref}
\setlength{\emergencystretch}{2em}

\newcommand{\G}{\mathbb{G}}
\newcommand{\GT}{\mathbb{G}_T}
\newcommand{\Zp}{\mathbb{Z}_p}
\newcommand{\Zq}{\mathbb{Z}_q}
\newcommand{\Fp}{\mathbb{F}_p}
\newcommand{\Setup}{\mathsf{Setup}}
\newcommand{\KeyGen}{\mathsf{KeyGen}}
\newcommand{\Store}{\mathsf{Store}}
\newcommand{\ChalGen}{\mathsf{ChalGen}}
\newcommand{\ProofGen}{\mathsf{ProofGen}}
\newcommand{\Verify}{\mathsf{Verify}}
\newcommand{\TagGen}{\mathsf{TagGen}}
\newcommand{\TagVerify}{\mathsf{TagVerify}}
\newcommand{\Response}{\mathsf{Response}}
\newcommand{\ProofVerify}{\mathsf{ProofVerify}}
\newcommand{\IndexGen}{\mathsf{IndexGen}}
\newcommand{\AuthGen}{\mathsf{AuthGen}}
\newcommand{\TrapdoorGen}{\mathsf{TrapdoorGen}}

\title{Extracted PDP Protocols from Reference Schemes}
\author{}
\date{}

\begin{document}
\maketitle

\section{Scope and Unified Reading}

This note extracts the PDP protocol cores from the four reference PDFs in
\texttt{refs}.  The scheme labels follow the naming style used by
\texttt{paper/FoldAudit.tex}; the leading ``20'' is removed from the year in
each label.

\begin{center}
\small
\begin{tabular}{p{0.28\textwidth}p{0.16\textwidth}p{0.46\textwidth}}
\toprule
PDF & Scheme label & PDP protocol core \\
\midrule
\texttt{ZhangTPDS2023.pdf} & ZhangTPDS23 & PCS-MPAP batch auditing over multiple files \\
\texttt{YuTC2025.pdf} & YuTC25 & HHF, polynomial commitment, and Tag-IMHT auditing \\
\texttt{XuTIFS2026.pdf} & XuTIFS26 & MPA, polynomial commitments plus Merkle authentication \\
\texttt{MiaoSCIS2026.pdf} & MiaoSCIS26 & multi-keyword searchable PDP \\
\bottomrule
\end{tabular}
\end{center}

The notation below preserves each paper's protocol structure.  When comparing
with FoldAudit, \(m\) may denote a file batch, \(n\) the number of blocks, \(s\)
the number of sectors per block, and \(c\) the number of challenged blocks.
Multi-server and multi-copy components in the source papers are normalized away
by taking a single storage server.  The extracted formulas therefore keep
multi-file auditing behavior while omitting cross-server aggregation and
diagnosis logic.

\section{ZhangTPDS23: PCS-MPAP Batch Auditing}

\subsection{Setup and Storage}

The system manager samples bilinear-group parameters
\((\G_1,\G_2,p,g,e)\), hash functions \(H:\{0,1\}^\ast\to\G_1\) and
\(H':\G_1\to\Zp\), and publishes
\[
    \mathsf{Para}=(\G_1,\G_2,g,e,H,H').
\]
The user samples \(x,\alpha\in\Zp^\ast\) and sets
\[
    sk=(x,\alpha),\qquad
    pk=(v=g^x,u=g^{\alpha x},g,g^\alpha,\ldots,g^{\alpha^{s-1}}).
\]
For each outsourced file \(F_l\), block \(j\) is encoded as
\[
    f_{l,j}(X)=\sum_{k=1}^{s} m_{l,j,k}X^{k-1}.
\]
The block authenticator is
\[
    \sigma_{l,j}
    =
    \left(H(Fid_l\Vert j)\cdot g^{f_{l,j}(\alpha)}\right)^x .
\]
The storage server checks a received tag by
\[
    e(\sigma_{l,j},g)
    \stackrel{?}{=}
    e\!\left(H(Fid_l\Vert j)\cdot g^{f_{l,j}(\alpha)},v\right).
\]

\subsection{Efficient Batch Auditing}

For a batch of \(d\) files, the smart contract samples
\[
    chal=(\{Q_l\}_{l=1}^{d},\{r_l\}_{l=1}^{d},\gamma,z),
\]
where \(Q_l=\{(a_j,v_j)\}\) is the challenge set for file \(l\), \(r_l\) is its
evaluation point, and \(\gamma,z\in\Zp\).

For the single storage server, the multi-file aggregated tag is
\[
    \sigma'
    =
    \prod_{l=1}^{d}
    \left(
        \prod_{(a_j,v_j)\in Q_l}\sigma_{l,a_j}^{v_j}
    \right)^{\beta_l},
    \qquad
    \beta_l=\gamma^{l-1}Z_{T\setminus\{r_l\}}(z),
\]
where \(T=\{r_1,\ldots,r_d\}\) and \(Z_S(X)=\prod_{r\in S}(X-r)\).  For each
file,
\[
    F_l(X)=\sum_{(a_j,v_j)\in Q_l} v_j f_{l,a_j}(X),
    \qquad
    S_l=F_l(r_l).
\]
The proof polynomials are
\[
\begin{aligned}
    F(X)
    &=
    \sum_{l=1}^{d}\gamma^{l-1}Z_{T\setminus\{r_l\}}(X)
    \bigl(F_l(X)-F_l(r_l)\bigr),\\
    H_b(X)&=\frac{F(X)}{Z_T(X)},\\
    L(X)
    &=
    \sum_{l=1}^{d}\beta_l
    \bigl(F_l(X)-F_l(r_l)\bigr)
    -Z_T(z)H_b(X).
\end{aligned}
\]
The server returns
\[
    W=g^{H_b(\alpha)},\qquad
    W'=g^{L(\alpha)/(\alpha-z)},\qquad
    E=\sum_{l=1}^{d}\beta_l S_l,
\]
and publishes
\[
    \mathsf{Proof}=(\sigma',W,W',E).
\]
The verifier checks
\[
    e(\sigma',g)\cdot e(\psi,v)
    \stackrel{?}{=}
    e(\zeta,v)\cdot e(W',u\cdot v^{-z}),
\]
where
\[
    \psi=g^{-E}W^{-Z_T(z)},\qquad
    \zeta=
    \prod_{l=1}^{d}
    \left(
        \prod_{(a_j,v_j)\in Q_l}H(Fid_l\Vert j)^{v_j}
    \right)^{\beta_l}.
\]

\section{YuTC25: HHF, Polynomial Commitment, and Tag-IMHT}

\subsection{Basic Auditing Protocol}

The device selects a secret \(\psi\) and publishes polynomial-commitment
parameters
\[
    \Psi_i=g^{\psi^i},\qquad i\in[0,n].
\]
For a chunk \(F_i=(F_{i,0},\ldots,F_{i,n-1})\), define
\[
    P_{F_i}(x)=F_{i,0}+F_{i,1}x+\cdots+F_{i,n-1}x^{n-1}.
\]
The device computes the homomorphic tag
\[
    \varpi_i=\mathsf{HHF}(P_{F_i}(\psi)),
\]
then commits to tags in a Tag-IMHT by hashing \(hc_i=C(\varpi_i)\) and
publishing the Merkle root \(h_r\).

Given challenge parameters \(cp=(k_1,k_2,c)\), the server derives
\[
    S=\{(id,a_{id})\},\qquad id=\pi_{k_1}(l),\quad a_{id}=f_{k_2}(l),
    \quad l\in[1,c],
\]
and \(z=f_{k_2}(c+1)\).  The proof polynomial is
\[
    P_{\mathsf{prf}}(x)=
    \sum_{id\in S} a_{id}P_{F_{id}}(x),
\]
with quotient
\[
    Z_{\mathsf{prf}}(x)=
    \frac{P_{\mathsf{prf}}(x)-P_{\mathsf{prf}}(z)}{x-z}.
\]
The server returns
\[
    \mathsf{prf}=\{Z_{\mathsf{prf}}(x),\{\varpi_i,\Omega_i\}_{i\in S},
    P_{\mathsf{prf}}(z)\},
\]
where \(\Omega_i\) is the auxiliary authentication information for the
Tag-IMHT path.

The verifier reconstructs \(h'_r\) from \(\{\varpi_i,\Omega_i\}\), checks
\(h'_r=h_r\), and computes
\[
    H_Z=\mathsf{HHF}(Z_{\mathsf{prf}}(\psi))
    =\prod_{i=0}^{n-1}\Psi_i^{\rho_i},
    \qquad
    H_{\mathsf{prf}}=H_Z^\psi
    =\prod_{i=1}^{n}\Psi_i^{\rho_{i-1}},
\]
where \(Z_{\mathsf{prf}}(x)=\sum_{i=0}^{n-1}\rho_i x^i\).  Verification checks
\[
    \prod_{i\in S}\varpi_i^{a_i}
    \stackrel{?}{=}
    H_{\mathsf{prf}}\cdot H_Z^{-z}\cdot
    \mathsf{HHF}(P_{\mathsf{prf}}(z)).
\]

\subsection{Privacy-Enhanced Verification}

The privacy variant changes the proof polynomial.  The server samples
\(\beta\in\Zp^\ast\), sets
\[
    B=\mathsf{HHF}(\beta),\qquad \eta=H(B),
\]
and constructs
\[
    P_{\mathsf{prf}}(x)=
    \sum_{id\in S}a_{id}P_{F_{id}}(x)+\beta\eta.
\]
It then returns
\[
    \mathsf{prf}=
    \{Z_{\mathsf{prf}}(x),\{\varpi_i,\Omega_i\}_{i\in S},
    P_{\mathsf{prf}}(z),B\}.
\]
After the Tag-IMHT check, the verifier accepts iff
\[
    B^\eta\cdot\prod_{i\in S}\varpi_i^{a_i}
    \stackrel{?}{=}
    H_{\mathsf{prf}}\cdot H_Z^{-z}\cdot
    \mathsf{HHF}(P_{\mathsf{prf}}(z)).
\]
The mask is a constant term; therefore it affects \(P_{\mathsf{prf}}(z)\) but
cancels out in the quotient \(Z_{\mathsf{prf}}(x)\).

\section{XuTIFS26: MPA Auditing}

\subsection{Setup, Storage, and Tag Verification}

A trusted center publishes
\[
    pp=(\G,\GT,q,g,\hat e,H).
\]
The user samples \(\alpha\in\Zq^\ast\) and publishes
\[
    pk=(g,g^\alpha,g^{\alpha^2},\ldots,g^{\alpha^{s-1}}).
\]
A file \(F_l\) is split into \(n\) blocks with \(s\) sectors:
\[
    F_l=\{m_{l,j,k}\}_{1\le j\le n,\,1\le k\le s}.
\]
Each block \(m_{l,j}\) is encoded as
\[
    f_{l,j}(X)=\sum_{k=1}^{s}m_{l,j,k}X^{k-1},
\]
and tagged by the KZG-style commitment
\[
    \sigma_{l,j}=g^{f_{l,j}(\alpha)}.
\]
The tags are hashed as Merkle leaves \(leaf_{l,j}=H(\sigma_{l,j})\), and the
root \(Root_l\) is published on-chain.

To validate received storage, the server recomputes
\[
    C_{l,j}=\prod_{k=1}^{s}(g^{\alpha^{k-1}})^{m_{l,j,k}}
    =g^{f_{l,j}(\alpha)}
\]
and accepts only if the Merkle root built from \(H(C_{l,j})\) matches
\(Root_l\).

\subsection{Challenge, Response, and Verification}

The blockchain commit-and-reveal challenge returns
\[
    Chal=(\{a_j,v_j\}_{j=1}^{c},r),
\]
where \(a_j\in[1,n]\), \(v_j\in\Zq^\ast\), and \(r\in\Zq^\ast\).

For file \(F_l\), the server computes
\[
    f_{M_l}(X)=\sum_{j=1}^{c}v_j f_{l,a_j}(X),\qquad
    val_l=f_{M_l}(r),
\]
\[
    h_l(X)=\frac{f_{M_l}(X)-val_l}{X-r},\qquad
    w_l=g^{h_l(\alpha)}.
\]
Equivalently, if \(h_l(X)=\sum_k h_l[k]X^k\), then
\[
    w_l=\prod_k(g^{\alpha^k})^{h_l[k]}.
\]
The per-file proof is
\[
    P_l=(val_l,w_l,\{\sigma_{l,a_j},MerklePath_{l,a_j}\}_{j=1}^{c}).
\]

For a batch of \(T\) files on a single server, the verifier aggregates
\[
    \sigma_{\mathsf{global}}
    =\prod_{l=1}^{T}\prod_{j=1}^{c}\sigma_{l,a_j}^{v_j},
    \qquad
    val_{\mathsf{global}}=\sum_{l=1}^{T}val_l,
    \qquad
    w_{\mathsf{global}}=\prod_{l=1}^{T}w_l.
\]
The multi-file KZG check is
\[
    \hat e(\sigma_{\mathsf{global}}/g^{val_{\mathsf{global}}},g)
    \stackrel{?}{=}
    \hat e(w_{\mathsf{global}},g^{\alpha-r}).
\]
The verifier also checks the supplied Merkle path for every challenged block
against the corresponding \(Root_l\).  Multi-copy and multi-provider fallback
diagnosis from the original MPA setting is omitted under the single-server
comparison model.

\section{MiaoSCIS26: Multi-Keyword Searchable PDP}

\subsection{Initialization, Indexes, and Authenticators}

The protocol starts from \(v\) files and keyword set
\[
    W=\{w_1,\ldots,w_K\}.
\]
The data owner initializes an \(a\times a\) matrix \(I_{a\times a}\), where
\(a\gg v\) and \(a\gg K\), to hide the true numbers of files and keywords.
Rows encode keyword membership vectors and columns encode files.

The data owner publishes bilinear parameters
\[
    param=\{\G_1,\GT,p,g,u,y,pk,e,H_1,H_2,H_3,h_1,
    f_{key},\psi_{key},\pi_{key},\mathsf{ssg}_{sk}\},
\]
where \(x\in\Zq^\ast\) is the secret key and \(y=g^x\).  For file \(F_i\)
with keyword set \(W_{F_i}=\{w_{b_1},\ldots,w_{b_{t_i}}\}\), define
\[
    H_{F_i}=\prod_{l=1}^{t_i}H_1(w_{b_l}).
\]
The match index is
\[
    \Omega_i=(H_{F_i}\cdot H_2(FID_i))^{-x}.
\]
The blinded row address for keyword \(w_k\) is \(\pi_o(w_k)\), and the
encrypted row is
\[
    ev_{\pi_o(w_k)}=v_{w_k}\oplus f_l(\pi_o(w_k)).
\]

For data blocks \(m_{i,j}\), the data owner creates index-switch values
\(\theta_{i,j}=(j,au_{i,j})\) and authenticators
\[
    \sigma_{i,j}
    =
    \left(
        H_{F_i}\cdot H_2(FID_i)\cdot H_3(au_{i,j})\cdot u^{m_{i,j}}
    \right)^x.
\]
The cloud can verify the uploaded authenticators by checking the aggregate
pairing relation
\[
    e\!\left(\prod_{i=1}^{v}\prod_{j=1}^{n}\sigma_{i,j},g\right)
    \stackrel{?}{=}
    e\!\left(
        \prod_{i=1}^{v}\prod_{j=1}^{n}
        H_{F_i}H_2(FID_i)H_3(au_{i,j})u^{m_{i,j}},
        y
    \right).
\]

\subsection{Trapdoor, Challenge, Proof Generation, and Verification}

For target keywords \(W'=\{w_1,\ldots,w_h\}\), the data owner sends
\[
    T_w=\{(\pi_o(w_i),f_l(\pi_o(w_i)))\}_{i=1}^{h}
\]
to the auditor.  The auditor samples \(k_1,k_2\in\Zq^\ast\) and challenge
count \(c\), and sends
\[
    Chal=(T_w,c,k_1,k_2)
\]
to the cloud.

Both auditor and cloud derive challenged block positions
\[
    Q=\{j=\psi_{k_1}(\beta)\}_{\beta=1}^{c}
\]
and per-file coefficients \(v_{i,j}\) from \(\pi_{k_2}\).  Using the trapdoor,
they decrypt matching index rows and derive the set of files that contain the
queried keywords.  Let \(R\) denote the set found by the cloud and \(T\) the set
found by the auditor.

The cloud computes the aggregate data value and aggregate authenticator
\[
    \mu=\sum_{i\in R}\sum_{j\in Q} e_{i,j}m_{i,j},
    \qquad
    T_{auth}=\prod_{i\in R}\prod_{j\in Q}\sigma_{i,j}^{e_{i,j}}.
\]
It samples \(r\in\Zq^\ast\), sets \(R_g=g^r\), and masks the aggregate data as
\[
    \mu_1=\mu-rh_1(T_w,R_g).
\]
The proof is
\[
    \mathsf{Proof}=(T_{auth},\mu_1,R_g).
\]

The auditor retrieves the match indexes \(\Omega_i\) and authenticator indexes
\(au_{i,j}\) for \(i\in T\) and \(j\in Q\).  It accepts iff
\[
e\!\left(
    T_{auth}\cdot
    \prod_{i\in T}\Omega_i^{\sum_{j\in Q}v_{i,j}},
    g
\right)
\stackrel{?}{=}
e\!\left(
    \left(\prod_{i\in T}\prod_{j\in Q}H_3(au_{i,j})^{v_{i,j}}\right)
    \cdot u^{\mu_1}\cdot R_g^{h_1(T_w,R_g)},
    y
\right).
\]
This equation simultaneously checks the matched-file set, challenged
authenticators, and the masked aggregate data value.

\subsection{File Dynamics}

To add a file \(F_z\), the data owner inserts a new keyword column in
\(I_{a\times a}\), recomputes the encrypted affected rows, calculates
\[
    H_{F_z}=\prod_{w\in W_{F_z}}H_1(w),\qquad
    \Omega_z=(H_{F_z}H_2(FID_z))^{-x},
\]
and generates
\[
    \sigma_{z,j}=
    \left(H_{F_z}H_2(FID_z)H_3(au_{z,j})u^{m_{z,j}}\right)^x.
\]
The blockchain receives \(\{add,FID_z,S_z,EI',\Omega_z\}\), while the cloud
receives \(\{\phi_z,\{m_{z,j}\}\}\) and verifies the uploaded authenticators.

To delete file \(F_z\), the data owner zeros the corresponding matrix column,
reencrypts the affected rows to form \(EI'\), publishes \(\{EI'\}\) on-chain,
and sends
\[
    \{delete,FID_z,\mathsf{ssg}_{sk}(delete\Vert FID_z)\}
\]
to the cloud for authenticated deletion.

\section{Comparison Pointers for FoldAudit}

The four extracted protocols motivate the comparison dimensions used in the
FoldAudit paper.  ZhangTPDS23 has distinct per-file evaluation
points in the PCS-MPAP batch proof but exposes an unmasked linear combination
of challenged evaluations.  YuTC25 masks the proof polynomial by adding a
constant term, but the quotient polynomial still reveals the non-constant part
of the challenged aggregate.  XuTIFS26 supports multi-file aggregation and
Merkle-authenticated updates, while its global check uses a shared evaluation
point.  MiaoSCIS26 provides keyword-filtered multi-file PDP with data masking,
but it is not a polynomial-opening batch protocol.

\end{document}
